\clearpage
\chap{Regularization selects a protected correction}\label{sec:regularized-repair}
A fixed budget asks how small the residual can become without exceeding an allowed correction size. Quadratic regularization reverses the emphasis. It asks for the best compromise between the residual and the size of the change. The resulting problem has a unique answer even when the exact protected equation has no solution.

That last sentence is both its usefulness and its limitation. The existence of a regularized answer is not a proof of exact realizability. The limiting behavior, rather than the mere existence of the minimizer, supplies the diagnosis.

\subsection{The normal equation is derived in the protected space}
For $\alpha>0$, define on $H=\ker p$
\[
 J_\alpha(a)=\norm{Ta-r}^2+\alpha\norm a^2.
\]
Both norms are the declared Hilbert norms. All minimizers below are corrections in $H$, so preservation is exact for every parameter, not just in the limit. This is ordinary quadratic Tikhonov regularization applied to the actual protected correction operator. The standard regularization background is discussed in \cite{clason-inverse}. The proof here does not require $T$ to be compact.

\begin{theorem}[A unique regularized correction]\label{WR06}
For every $\alpha>0$, there is a unique minimizer
\[
 a_\alpha=(T^*T+\alpha I)^{-1}T^*r.
\]
It satisfies $T^*(r-Ta_\alpha)=\alpha a_\alpha$. For every $a\in H$,
\begin{equation}\label{eq:regularization-energy}
 J_\alpha(a)-J_\alpha(a_\alpha)
 =\norm{T(a-a_\alpha)}^2+\alpha\norm{a-a_\alpha}^2.
\end{equation}
\end{theorem}
\begin{proof}
Let $B=T^*T+\alpha I$. Then $B$ is bounded and self-adjoint, and
\[
 \langle Ba,a\rangle=\norm{Ta}^2+\alpha\norm a^2
 \ge\alpha\norm a^2.
\]
Cauchy--Schwarz gives $\norm{Ba}\ge\alpha\norm a$. Thus $B$ is injective and has closed range: a convergent sequence $Ba_n$ forces $(a_n)$ to be Cauchy, and continuity identifies the limit. If $y$ is orthogonal to the range, self-adjointness gives $By=0$, hence $y=0$. The closed range is therefore dense and equals $H$. Its inverse exists and has norm at most $1/\alpha$.

Define $a_\alpha=B^{-1}T^*r$. Expanding $J_\alpha(a_\alpha+h)$, the two linear terms cancel by the normal equation, leaving \eqref{eq:regularization-energy}. The right side is nonnegative and can vanish only for $h=0$. This proves existence, global minimality, and uniqueness.
\end{proof}
\written{WR06}

A declared regularization parameter is part of the problem. A calculation that changes $\alpha$ is solving a different optimization problem, even if the protected data and forward operator stay fixed. On the other hand, the preservation condition itself does not change with $\alpha$.

\subsection{Coordinate bases must retain the correction metric}
In finite dimensions, let $B$ have linearly independent columns spanning $\ker p$. Write $a=Bc$. The coefficient vector $c$ is a representation of the correction, not necessarily an isometric one. Set
\[
 S=dB,\qquad G=B^{\mathsf T}B.
\]
The correct coordinate equation is
\begin{equation}\label{eq:gram-regularization}
 (S^{\mathsf T}S+\alpha G)c=S^{\mathsf T}r.
\end{equation}
The matrix $G$ is positive definite because the columns are independent. Replacing $G$ by the identity changes the meaning of correction size unless the chosen basis is orthonormal.

For example, let $d(a_1,a_2)=a_1$, $p(a_1,a_2)=a_1-a_2$, $r=1$, and $\alpha=1$. A basis of $\ker p$ is the column $B=(1,1)^{\mathsf T}$. Here $S=1$ and $G=2$, so the correction is $(1/3,1/3)$. Penalizing the coordinate with $G=1$ instead gives $(1/2,1/2)$. The latter is not the optimizer for the ambient Euclidean norm. The difference is a metric change, not a discrepancy between equivalent calculations.

The executable example in this edition constructs the nullspace basis over the rationals, retains its Gram matrix, solves \eqref{eq:gram-regularization}, and returns both the correction and an exact normal-equation residual. It also handles a zero-dimensional protected kernel without pretending that a nonzero correction exists.

\subsection{Parameter monotonicity and the fixed-budget problem}
\begin{proposition}[Regularized answers lie on the budget frontier]\label{WR07}
If $0<\alpha<\gamma$, then
\[
 \norm{a_\alpha}\ge\norm{a_\gamma},\qquad
 \norm{r-Ta_\alpha}\le\norm{r-Ta_\gamma}.
\]
At $M_\alpha=\norm{a_\alpha}$,
\[
 \beta(M_\alpha)=\norm{r-Ta_\alpha}.
\]
When $v_\alpha=r-Ta_\alpha\ne0$, its normalization is an optimizing witness in WR04 at that budget.
\end{proposition}
\begin{proof}
Write $A_\alpha=\norm{a_\alpha}^2$ and $R_\alpha=\norm{r-Ta_\alpha}^2$. Optimality at the two parameters gives
\[
 R_\alpha+\alpha A_\alpha\le R_\gamma+\alpha A_\gamma,
 \qquad
 R_\gamma+\gamma A_\gamma\le R_\alpha+\gamma A_\alpha.
\]
Adding proves $(\gamma-\alpha)(A_\alpha-A_\gamma)\ge0$, and substitution into the first inequality gives $R_\alpha\le R_\gamma$.

If a correction with norm at most $M_\alpha$ had a strictly smaller residual, its $J_\alpha$ value would also be smaller. Thus $a_\alpha$ attains the budget minimum. The normal equation gives
\[
 T^*v_\alpha=\alpha a_\alpha,\qquad
 \langle r,v_\alpha\rangle-\alpha\norm{a_\alpha}^2
 =\norm{v_\alpha}^2.
\]
Divide the second equation by $\norm{v_\alpha}$ to obtain equality in the dual formula.
\end{proof}
\written{WR07}

This links a penalized calculation to a hard-budget certificate. It does not say that every larger budget must produce a different answer. Once the least-squares floor is attained, excess budget can be irrelevant. If $T^*r=0$, every regularized correction is zero and the entire profile is constant.

\subsection{The limit separates projection from actual realizability}
Keep $C=\overline{\ran T}$, $b=P_Cr$, and $q=\norm{r-b}$. Since $r-b\in\ker T^*$, the regularized corrections depend on $b$, not on the component that is orthogonal to the closure of all attainable residuals.

\begin{theorem}[The regularization limit alternative]\label{WR08}
As $\alpha\downarrow0$,
\[
 Ta_\alpha\longrightarrow b,\qquad
 \norm{r-Ta_\alpha}\longrightarrow q,\qquad
 \alpha\norm{a_\alpha}^2\longrightarrow0.
\]
If $b\in\ran T$, let $a^\dagger$ be its unique minimum norm preimage. Then
\[
 a_\alpha\longrightarrow a^\dagger\quad\text{in norm}.
\]
If $b\notin\ran T$, then $\norm{a_\alpha}\to\infty$.
\end{theorem}
\begin{proof}
For any fixed $a\in H$, minimality gives
\[
 q^2\le\norm{r-Ta_\alpha}^2
 \le J_\alpha(a_\alpha)\le\norm{r-Ta}^2+\alpha\norm a^2.
\]
Choose $a$ with residual square arbitrarily close to $q^2$, then let $\alpha\downarrow0$. It follows that both the residual square and the minimum objective tend to $q^2$. Their difference is $\alpha\norm{a_\alpha}^2$, so that term tends to zero. The orthogonal identity \eqref{eq:range-decomposition} then yields $Ta_\alpha\to b$.

If $b\in\ran T$, its preimage set is a nonempty closed affine space. Projection of zero gives $a^\dagger$. Comparison with $a^\dagger$ yields
\[
 \norm{Ta_\alpha-b}^2+\alpha\norm{a_\alpha}^2
 \le\alpha\norm{a^\dagger}^2,
\]
so $\norm{a_\alpha}\le\norm{a^\dagger}$. Along any sequence of parameters tending to zero, a weakly convergent subsequence has an exact preimage of $b$ as its limit. The norm bound and weak lower semicontinuity make that limit the minimum norm preimage $a^\dagger$. Along that subsequence, the norms converge to $\norm{a^\dagger}$, and weak convergence together with convergence of norms gives strong convergence. If the whole parameter family failed to converge strongly, a sequence staying a fixed distance away would have just such a strongly convergent subsequence, a contradiction.

If $b\notin\ran T$ and the correction norms failed to tend to infinity, some parameter sequence tending to zero would have bounded corrections. Their images tend to $b$, contradicting WR01.
\end{proof}
\written{WR08}

The vanishing penalty $\alpha\norm{a_\alpha}^2$ is not a boundedness certificate for $a_\alpha$. Its coefficient is tending to zero. In the escaping case, the penalty and residual can both tend to zero while the actual corrections diverge.

\begin{center}
\begin{tabularx}{\linewidth}{@{}p{.17\linewidth}p{.26\linewidth}X@{}}
\toprule Residual floor & Is $b$ attained? & Limit of the protected corrections\\\midrule
$q=0$ & Yes & Strong convergence to the nearest exact protected correction.\\[3pt]
$q=0$ & No & Residual vanishes, correction norms diverge, no exact protected answer.\\[3pt]
$q>0$ & Yes & Strong convergence to the least norm least-squares correction, not an exact answer.\\[3pt]
$q>0$ & No & A positive residual floor and diverging correction norms occur together.\\\bottomrule
\end{tabularx}
\end{center}

The last case is not hypothetical. Let $Tu=(u_n/n)$ on $\ell^2$, define $Su=(Tu,0)$ into $\ell^2\oplus\R$, and take $r=((1/n),1)$. Then $b=((1/n),0)$ is in the closure but not the range, while $q=1$. The same escaping corrections as before coexist with an irreducible orthogonal residual. Positive floor and unbounded correction cost are independent features.

\subsection{Noise is controlled at a fixed parameter}
\fixedplate{34_noise_filter}{For one singular value $\sigma$, the response to residual-data noise is $\sigma/(\sigma^2+\alpha)$. Its maximum occurs at $\sigma=\sqrt\alpha$ and equals $1/(2\sqrt\alpha)$. The plot evaluates this formula and the universal horizontal bounds at three explicit positive parameters.}{fig:noise-filter}

\begin{theorem}[Data sensitivity and a convergent parameter choice]\label{WR09}
For two residual data $r,\widetilde r$ with $\norm{r-\widetilde r}\le\delta$ and the same $T,\alpha$,
\[
 \norm{a_\alpha(r)-a_\alpha(\widetilde r)}
 \le\frac{\delta}{2\sqrt\alpha},\qquad
 \norm{T(a_\alpha(r)-a_\alpha(\widetilde r))}\le\delta.
\]
The coefficient $1/(2\sqrt\alpha)$ is sharp over bounded linear Hilbert operators. If $b=P_Cr\in\ran T$, $\alpha(\delta)\to0$, and $\delta/\sqrt{\alpha(\delta)}\to0$, then the noisy corrections converge strongly to $a^\dagger$ as $\delta\to0$.
\end{theorem}
\begin{proof}
Let $h=a_\alpha(\widetilde r)-a_\alpha(r)$ and $e=\widetilde r-r$. Subtracting the normal equations and pairing with $h$ gives
\[
 \norm{Th}^2+\alpha\norm h^2=\langle e,Th\rangle
 \le\delta\norm{Th}.
\]
Put $t=\norm{Th}$. Then $t\le\delta$ and
\[
 \alpha\norm h^2\le\delta t-t^2
 =\frac{\delta^2}{4}-(t-\delta/2)^2\le\frac{\delta^2}{4}.
\]
This proves the two estimates. For the scalar operator $T=\sqrt\alpha\,I$, the solution map is multiplication by $1/(2\sqrt\alpha)$, which proves sharpness. Finally, split the noisy error into the data perturbation at the chosen parameter and the exact-data regularization error. The first tends to zero by the displayed estimate and the second by WR08.
\end{proof}
\written{WR09}

The theorem does not manufacture an exact solution when $r$ is outside the range. For the positive-floor attainable case it converges to a least-squares correction. For the unattainable projected datum it cannot produce a bounded exact-data limit. The range hypothesis must therefore precede the parameter-choice conclusion.


\begin{insight}
Regularization gives a unique, exactly protected answer at every positive parameter. Whether those answers converge to an exact realization, converge only to a best approximation, or escape is determined by the projected datum's actual range membership. A smooth regularized path does not erase that distinction.
\end{insight}
